\section{CA Implementation Extensions}
\label{app:dgm-extensions}

This appendix collects CA mechanics and large-model details that are useful for implementations but secondary to the main refine-or-repair argument. The main experiments instantiate the fixed-bank CA layer from Section~\ref{sec:main-ca}. The graph-growing material below explains how the same refine-or-repair principle can be used when concepts and distinction edges are created online.

\subsection{Trees as a hard-routing special case}

A tree stores distinctions as a hierarchy. A mistake near the root changes every downstream routing decision. A graph stores distinctions locally between concepts. Creating a new concept does not rewrite old paths; it adds a local alternative and edges explaining why the new concept differs from its neighbors. Trees are still contained in the view below, but as a constrained hard-routing case.

\begin{proposition}[Trees as generalized hard-CAs, informal version of Proposition~\ref{prop:tree-as-dgm:formal}]
\label{prop:tree-as-dgm:informal}
Given a binary distinction tree, construct one generalized hard-route CA concept for each leaf and attach the path distinctions as hard eligibility constraints for that concept. If the generalized hard route selects the unique eligible leaf concept and each concept memory equals the tree leaf predictor, then the generalized CA and the tree make identical predictions.
\end{proposition}

Thus the distinction-graph view contains trees after allowing arbitrary concept-level eligibility constraints. A decision tree is the sparse case in which distinctions are forced into a rooted hierarchy and routing is hard and single-path.

\subsection{CA data-structure cost}

Let
\begin{align*}
Q(h)=\sum_{j\in N(h)}\abs{\operatorname{Inc}_t(j)}
\end{align*}
be the number of incident edge gates evaluated for the selected candidates. With exact nearest-neighbor search, prediction costs
\begin{align*}
O(N_td+Q(h)d+\abs{N(h)}d+C_{\mathrm{local}}),
\end{align*}
where $C_{\mathrm{local}}$ is the cost of evaluating the selected local memories and the $Q(h)d$ term evaluates dense linear edge gates. With approximate nearest-neighbor indexing, the $N_td$ term can be replaced by the query cost of the index \citep{indyk1998approximate,johnson2019billion}. Creating an anchor-centroid pair costs $O(d)$ plus the cost of initializing its memory and the bookkeeping used for anchor-order queries. Adding $q$ guarded distinction edges costs $O(qd)$. If each concept stores a sparse categorical memory with average support size $\bar s$, runtime memory is
\begin{align*}
O(N_t(2d+\bar s)+\abs{E_t}d).
\end{align*}
An embedding-mean variant instead uses $O(N_t(2d+d_e)+\abs{E_t}d)$ memory, but it is not a calibrated log-loss memory without an additional probability model.

\begin{theorem}[CA data-structure bound, informal version of Theorem~\ref{thm:dgm-complexity:formal}]
\label{thm:dgm-complexity:informal}
For exact search over $N_t$ centroids in $\R^d$, one edge-gated CA prediction costs $O(N_td+Q d+\abs{N}d+C_{\mathrm{local}})$, where $N=N(h)$ and $Q=Q(h)$ is an edge count. One observation that repairs a selected concept costs this plus the guarded local repair scoring cost. One observation that creates a new anchor-centroid pair and $q$ guarded distinction edges costs an additional $O(d+qd)$ plus anchor-order query bookkeeping. With sparse categorical memories of average support $\bar s$, the runtime state uses $O(N_t(2d+\bar s)+\abs{E_t}d)$ memory.
\end{theorem}

\subsection{Ridge local memory}

For few-shot adaptation or local readout repair, a memory stores support embeddings $X_{\mathrm{sup}}\in\R^{n\times d}$ and labels $Y\in\R^{n\times C}$. The closed-form adapter is the dual ridge/Tikhonov solution \citep{tikhonov1977solutions,hoerl1970ridge}:
\begin{align*}
\Gamma=(X_{\mathrm{sup}}X_{\mathrm{sup}}^\top+\lambda I)^{-1}Y,
\qquad
\operatorname{logits}(z)=zX_{\mathrm{sup}}^\top\Gamma.
\end{align*}
Projection Memory and CPM are advanced repair rules inside $M_j$ after the graph has fixed the relevant local concept; Appendix~\ref{sec:projection-bridge} treats that relation directly.

\subsection{Training modes for representation drift}

Anchors, centroids, and edges live in query space. If the encoder or query map changes rapidly while the CA state persists, old anchors and edges describe stale geometry. Reports should therefore state which training mode is used.

\paragraph{Mode A: frozen encoder.}
Train the backbone and query map first, then freeze $\phi_\theta$, $Q$, and the backbone distribution used for CA smoothing while running CA runtime learning. This is the cleanest protocol for the first experiments and is the default interpretation of the theory.

\paragraph{Mode B: EMA query encoder.}
Maintain a slowly moving query encoder
\begin{align*}
\bar\theta\leftarrow \beta\bar\theta+(1-\beta)\theta,
\end{align*}
and create anchors, centroids, and edge generators in the query space of $\phi_{\bar\theta}$ and $Q_{\bar\theta}$. The backbone can still train, but the CA state is tied to the slower geometry.

\paragraph{Mode C: periodic state rebuild.}
During end-to-end training, rebuild or clear the CA state every $K$ optimization steps, or rebuild it from a replay buffer under the current query map. The formal guarantees then apply within a rebuild interval. Cross-interval state carryover is a systems approximation, not a claim of anchor-stable refinement.

\subsection{Non-canonical adapter extension}

Beyond the categorical next-token memory in Appendix~\ref{sec:large-model-training}, a richer optional adapter-style layer stores a local linear repair map, related to parameter-efficient adapter and LoRA modules but updated as a local CA memory \citep{houlsby2019parameter,hu2022lora}. It requires an explicit residual target and should not be conflated with the calibrated next-token memory mechanism:
\begin{align*}
\calS_t^{\mathrm{adp}}=\{a_j,c_j,W_j,n_j\}_{j=1}^{N_t}\cup G_t,
\qquad
W_j\in\R^{d\times r}.
\end{align*}
Let $K:\R^d\to\R^r$ be a learned or fixed key projection and set $k_t=Kh_t$. The adapter read is
\begin{align*}
\Delta h_t=\sum_{j\in N(q_t)}\alpha_j(q_t)W_jk_t,
\qquad
\tilde h_t=h_t+\gamma_\theta(h_t)\Delta h_t.
\end{align*}
Here $W_jk_t$ is a hidden-space residual, so the write rule requires an explicit residual target $v_t\in\R^d$. Let $P_{\mathrm{emb}\to h}:\R^{d_e}\to\R^d$ be a chosen projection from output-embedding space to hidden-residual space. Natural choices include a token-embedding target $P_{\mathrm{emb}\to h}e_{x_{t+1}}$, a teacher-forced output-error target $P_{\mathrm{emb}\to h}(e_{x_{t+1}}-\hat e_t)$, or a distillation target $h_t^{\mathrm{teacher}}-h_t$. Without such a target, the adapter is only an architectural extension.

Given a local residual target, the projected rank-one repair for a selected concept is the same residual-correction pattern that appears in passive-aggressive updates, RLS/ridge-style online least squares, and delta-rule fast-weight memories \citep{crammer2006online,vovk1997competitive,schlag2021linear}:
\begin{align*}
W_j^+=W_j+
\frac{\eta\alpha_j(q_t)}
{1+\eta\alpha_j(q_t)\norm{k_t}_2^2}
\big(v_t-W_jk_t\big)k_t^\top.
\end{align*}
Equivalently, a memory may maintain small ridge sufficient statistics and solve the local adapter in closed form \citep{tikhonov1977solutions,hoerl1970ridge}.

\subsection{Cost and systems notes}

During sequential detached training, backpropagation sees the selected read computation, while the CA state is updated in the same order as streaming inference. The important comparison is with test-time gradient adaptation \citep{sun2020test,wang2021tent,niu2022eata,zhang2022memo,sun2025learning}. A TTT-style layer typically performs an inner backward pass at inference time. A CA layer performs retrieval, sparse local read, and closed-form or sufficient-statistic repair. The cost is forward-only:
\begin{align*}
\text{CA inference}=\text{backbone forward}+\text{retrieval}+\text{sparse local repair}.
\end{align*}

\begin{theorem}[CA layer cost, informal version of Theorem~\ref{thm:dgm-layer-cost:formal}]
\label{thm:dgm-layer-cost:informal}
Consider a CA adapter extension with hidden dimension $d$, query dimension $d_q$, local rank $r$, candidate count $\abs{N}$, selected incident edge-gate count $Q_N$, and concept-bank query cost $C_{\mathrm{ret}}(N_t,d_q)$. For one token, the read overhead is
\begin{align*}
O(dd_q+C_{\mathrm{ret}}(N_t,d_q)+Q_N d_q+\abs{N}dr).
\end{align*}
If the write path is detached and uses rank-one local repair, the write overhead is
\begin{align*}
O(dd_q+C_{\mathrm{ret}}(N_t,d_q)+Q_N d_q+\abs{N}(d_q+dr)).
\end{align*}
For a block of length $L$, the CA state history need not be stored for backpropagation; the additional activation memory is the selected read computation, not the full sequence of memory states. The CALM mixer replaces the $\abs{N}dr$ adapter term by sparse concept lookup, edge-gated value aggregation, and detached graph-write costs.
\end{theorem}

The concept bank can be sharded by the same principle used in retrieval and sparse expert systems. Each device owns a subset of concepts. A query is routed to a small set of shards, each shard returns local top-$k$ candidates, and a final top-$k$ merge selects the memories used by the read path. Only selected memories receive gradients. Detached writes are sent back to the owning shards.

Budgeting is necessary for long training. A practical CA layer maintains $N_t\le N_{\max}$, merges only compatible unprotected concepts, and otherwise evicts or compresses low-score memories. These budget actions are maintenance operations, not learning operations in the formal refine-or-repair sense. The refine-or-repair theorems apply between maintenance events; merge, eviction, and compression may destroy information and must be evaluated as resource-management approximations.

\subsection{Closed-form local repair}

Ridge memory is the main closed-form adapter used in our checks. Let $X_{\mathrm{sup}}\in\R^{n\times d}$ be normalized support embeddings stored by a CA memory and let $Y\in\R^{n\times C}$ be one-hot consequences. The adapter computes the dual ridge/Tikhonov solution \citep{tikhonov1977solutions,hoerl1970ridge}
\begin{align*}
\Gamma=(X_{\mathrm{sup}}X_{\mathrm{sup}}^\top+\lambda I)^{-1}Y,
\qquad
\operatorname{logits}(z)=zX_{\mathrm{sup}}^\top\Gamma.
\end{align*}
This is projection-style repair: it fits the local adapter in one linear solve rather than by iterated backpropagation. In a large model, the same idea can be applied per concept, per document, or per retrieved shard after the relevant distinction structure has been fixed.
